\section{Finite diagnostics and conclusion}
\label{sec:computations-conclusion}

The theorems above are analytic and do not depend on computation.  We
nevertheless supply a deterministic finite diagnostic for three purposes:
to test the exact divisor split, to expose the slow finite-scale behavior of
the tail, and to illustrate the periodic-energy quantity in
\cref{thm:finite-conductor-energy-escape}.

\subsection{FFT defect certificate}
\label{subsec:FFT-certificate}

The script \texttt{experiments/defect\_shift\_diagnostics.py} constructs
$\mu$ and $\Lambda$ by sieving.  For each cutoff it forms $a_D$ by adding
$-\mu(d)\log d$ on the multiples of $d$, puts $r_D=\Lambda-a_D$, and computes
all correlations
\[
 \sum_{X<n\le2X}r_D(n)\Lambda(n+h),
 \qquad 1\le h\le H,
\]
with one zero-padded FFT convolution.  Direct dot products at selected
shifts and a second FFT for the short part verify
\begin{equation}\label{eq:finite-computational-identity}
 \cC_h(X)-\cC_{h,\le D}(X)-\cT_{h,>D}(X)=0
\end{equation}
up to floating-point roundoff.

The committed run uses $X=2^{18}$, $H=2^{15}$, and the exploratory cutoffs
$D=\operatorname{round}(X^\theta)$.  The even-shift summaries are shown in
\cref{tab:finite-defect-quantiles}.

\begin{table}[ht]
\centering
\caption{Finite even-shift defect diagnostics.  Each statistic is normalized
by $X$; it is not an asymptotic estimate.}
\label{tab:finite-defect-quantiles}
\begin{tabular}{@{}rrrrr@{}}
\toprule
$\theta$ & $D$ & median $|\cT|/X$ & $90\%$ quantile & RMS \\
\midrule
$0.25$ & $23$  & $0.3053$ & $1.8429$ & $0.9964$ \\
$0.35$ & $79$  & $0.2801$ & $0.9830$ & $0.5807$ \\
$0.45$ & $274$ & $0.1284$ & $0.4629$ & $0.3033$ \\
\bottomrule
\end{tabular}
\end{table}

The maximum discrepancy in \eqref{eq:finite-computational-identity} over
the three runs is below $5.8\times10^{-10}$.  The sizable normalized tails
in \cref{tab:finite-defect-quantiles} are a useful warning: the theorem
contains a non-explicit logarithmic cutoff constant, and small power cutoffs
at $X=2^{18}$ do not numerically realize its asymptotic regime.

\subsection{Periodic projection diagnostic}
\label{subsec:periodic-projection-diagnostic}

For every residue class modulo $q$, the script averages $\Lambda-1$ on that
class and evaluates the exact projection formula
\eqref{eq:projection-exact-formula}.  Selected captured-energy ratios appear
in \cref{tab:finite-periodic-energy}.

\begin{table}[ht]
\centering
\caption{Finite ratios
$\|P_q(\Lambda-1)\|_2^2/\|\Lambda-1\|_2^2$ on $(N,2N]$.}
\label{tab:finite-periodic-energy}
\begin{tabular}{@{}rrrr@{}}
\toprule
$N$ & $q=2$ & $q=30$ & $q=210$ \\
\midrule
$2^{15}$ & $0.1022$ & $0.2811$ & $0.3462$ \\
$2^{16}$ & $0.0957$ & $0.2632$ & $0.3237$ \\
$2^{17}$ & $0.0895$ & $0.2462$ & $0.3026$ \\
$2^{18}$ & $0.0844$ & $0.2322$ & $0.2852$ \\
\bottomrule
\end{tabular}
\end{table}

The decrease is qualitatively consistent with
\eqref{eq:captured-energy-ratio}, whose convergence is only logarithmic.
No statistical fit or extrapolation is used in a proof.

\begin{samepage}
The full machine-readable record is
\texttt{experiments/data/defect-shift-certificate.json}.  Its canonical
payload SHA-256 is
\begin{center}
\small\ttfamily
6F822312CE0699722A0830F98E91189823F3AB9ABE1F47810340E081AFFDAB30.
\end{center}
Six unit tests check the arithmetic arrays, the complete-cutoff identity,
the FFT convention, Euler's totient, the exact periodic projection, and
small-input/generator handling.
\end{samepage}

\begin{remark}[Computational boundary]
The finite values do not verify the hypothesis of an unproved conjecture,
choose the exceptional set in \cref{thm:complete-defect-almost-all}, or
provide evidence that $h=2$ is regular.  Their role is reproducibility and
finite algebraic validation only.
\end{remark}

\subsection{Conclusion and the remaining route}
\label{subsec:conclusion}

TPC-6 isolated the full large-source-divisor remainder and proved that its
cancellation at a fixed nonzero even shift is exactly the missing prime-pair
content.  The present paper replaces the fixed-shift quantifier in that
ledger by an almost-all-shift quantifier:
\begin{equation}\label{eq:conclusion-ledger}
\begin{gathered}
\text{uniform Bombieri--Vinogradov}\\[-2pt]
\text{short-divisor transfer}
\quad+\quad
\text{MRT almost-all prime-pair theorem}
\\
\Longrightarrow
\text{complete large-divisor defect is small at almost all shifts}.
\end{gathered}
\end{equation}
This gives an explicit estimate for that object within the TPC-6 ledger, but
only in shift density.  The fixed-coordinate spike in
\cref{prop:fixed-coordinate-spike} proves the quantifier gap.  The periodic
energy escape in \cref{thm:finite-conductor-energy-escape} and the
shifted-prime M\"obius slice \eqref{eq:outer-slice-shifted-primes} mark
boundaries for two natural strategies that might otherwise be proposed to
specialize \eqref{eq:conclusion-ledger} to a prescribed even shift.

One plausible route toward a subsequent fixed-shift advance would seek a
coefficient-stable nonlocal estimate.  A rigorous but narrower intermediate
target is a low-conductor second-coefficient transfer, in which a periodic
factor observable is passed through the prime target by
Bombieri--Vinogradov theory.  Such a theorem can describe exact arithmetic
involutions at growing polylogarithmic conductor, but
\cref{thm:finite-conductor-energy-escape} shows why it must not be confused
with approximation of the full weight $\Lambda-1$.  The broad fixed-shift
Type~II problem remains open.
